\chapter{Endotracheal Intubation}

\section*{Overview}
Endotracheal intubation involves inserting a flexible plastic tube (endotracheal tube) through the mouth into the trachea to maintain a patent airway and facilitate mechanical ventilation.

\section*{Alternative Names}
Intubation, Airway placement

\section*{Principle}
A laryngoscope blade lifts the tongue and epiglottis to expose the glottis, and an endotracheal tube is advanced through the vocal cords into the trachea. Inflating the cuff seals the airway, permitting positive-pressure ventilation and protecting the lungs from aspiration.
\section*{History}
William Macewen performed the first tactile oral intubation in 1878. Chevalier Jackson developed the laryngoscope in 1907, enabling direct visual intubation.

\section*{Why This Test or Procedure Is Ordered}
Ordered for patients with respiratory failure, airway obstruction, severe coma (Glasgow Coma Scale <= 8), or during general anesthesia.

\section*{Is This a Primary or Secondary Test or Procedure}
Primary definitive method for airway management and mechanical ventilation.

\section*{How This Test or Procedure Is Performed}
A laryngoscope is used to visualize the vocal cords. The endotracheal tube is passed through the vocal cords into the mid-trachea. The cuff at the end of the tube is inflated.

\section*{What Preparation Is Needed}
Pre-oxygenate the patient. Prepare suction, bag-valve mask, laryngoscope, tube, stylet, and medications (induction agent and paralytic).

\section*{Contraindications and Precautions}
None in emergency airway compromise. Relative contraindications include severe cervical spine injury or facial trauma (which may require a surgical airway).

\section*{Risks}
Esophageal intubation, hypoxia, dental injury, vocal cord injury, aspiration of gastric contents, or mainstem bronchus intubation.

\section*{What Happens After the Test or Procedure}
Confirm tube placement using end-tidal CO2 (capnography) and bilateral breath sounds. Secure the tube and obtain a chest X-ray to confirm tip position (mid-trachea).

\section*{Understanding Your Results}
Positive capnography waveform and equal breath sounds verify correct placement. Chest X-ray should show the tube tip 3-5 cm above the carina.

\section*{Normal Results}
Tube tip located in the mid-trachea; symmetric chest rise; bilateral breath sounds; end-tidal CO2 detection.

\section*{Follow-up Tests and Procedures}
Mechanical ventilation management, sedation titration, and daily evaluation for weaning and extubation.

\section*{References}
\begin{enumerate}
  \item MedlinePlus. Endotracheal Intubation. U.S. National Library of Medicine; 2025.
  \item Current Medical Diagnosis and Treatment. McGraw-Hill; 2025.
\end{enumerate}

\clearpage
